%==============================================================================
\section{Thuật toán Decision Tree (Cây quyết định)}
\label{sec:decision-tree}
%==============================================================================

\begin{dinhnghia}[title={Decision Tree Algorithm}]
Thuật toán cây quyết định là thuật toán phân cấp dạng \textbf{cây} dùng để phân loại hoặc dự đoán kết quả theo tập quy tắc.
Chia dữ liệu thành các tập con theo giá trị feature đầu vào; lặp đệ quy cho đến khi mỗi tập con thuộc cùng một lớp hoặc cùng giá trị biến mục tiêu.
Cây thu được là tập quy tắc quyết định để dự đoán hoặc phân loại dữ liệu mới.
\end{dinhnghia}

\begin{phuongphap}[title={Chọn feature tách tại mỗi nút}]
Tại mỗi nút, chọn feature \textbf{tốt nhất} để tách dữ liệu --- feature mang lại nhiều \textbf{information gain} nhất hoặc giảm \textbf{entropy} nhiều nhất.
Information gain đo lượng thông tin thu được khi tách tại một feature; entropy đo độ ngẫu nhiên/hỗn loạn trong dữ liệu.
Thuật toán dùng các đại lượng này để quyết định feature tách tại mỗi nút.
\end{phuongphap}

\begin{vidu}[title={Ví dụ cây nhị phân}]
Ví dụ cây nhị phân dự đoán người ``fit'' hay ``unfit'' theo tuổi, thói quen ăn uống và tập luyện:

\tsimg{05_decision_tree/img01.jpg}

Trong cây trên, các \textbf{câu hỏi} là nút quyết định (decision nodes); \textbf{kết quả cuối} là lá (leaves).
\end{vidu}

\subsection{Các loại thuật toán Decision Tree}

\begin{dinhnghia}[title={Classification Tree}]
Cây phân loại dùng để xếp dữ liệu vào các \textbf{lớp} hoặc \textbf{hạng mục} khác nhau.
Tách tập con theo feature và gán mỗi tập con một lớp.
\end{dinhnghia}

\begin{dinhnghia}[title={Regression Tree}]
Cây hồi quy dùng để dự đoán giá trị \textbf{số} (biến liên tục).
Tách tập con theo feature và gán mỗi tập con một giá trị số.
\end{dinhnghia}

\subsection{Triển khai thuật toán Decision Tree}

\subsubsection{Chỉ số Gini}

\begin{dinhnghia}[title={Gini Index}]
Gini Index là tên hàm \textbf{cost} đánh giá các tách nhị phân trong dataset, làm việc với biến mục tiêu phân loại Success/Failure (hoặc hai lớp).
Gini càng \textbf{cao} thì tập con càng \textbf{đồng nhất} (homogeneity).
Gini hoàn hảo = 0; tệ nhất = 0.5 (bài toán 2 lớp).
\end{dinhnghia}

\begin{phuongphap}[title={Tính Gini cho một tách}]
\begin{enumerate}
  \item Tính Gini cho từng nút con bằng công thức $p^2 + q^2$ (tổng bình phương xác suất thành công và thất bại).
  \item Tính Gini cho tách bằng \textbf{Gini có trọng số} của từng nút con.
\end{enumerate}
Thuật toán CART (Classification and Regression Tree) dùng phương pháp Gini để tạo tách nhị phân.
\end{phuongphap}

\subsubsection{Tạo tách (Split Creation)}

\begin{dinhnghia}[title={Split}]
Một tách gồm một \textbf{thuộc tính} trong dataset và một \textbf{giá trị}.
Tách dataset thành hai danh sách hàng (nhóm trái/phải) theo chỉ số thuộc tính và giá trị tách.
Sau khi có hai nhóm, tính giá trị tách bằng Gini ở bước trên; giá trị tách quyết định thuộc tính nằm ở nhóm nào.
\end{dinhnghia}

\begin{phuongphap}[title={Đánh giá mọi tách}]
Kiểm tra mọi giá trị gắn với mỗi thuộc tính như ứng viên tách.
Tìm tách tốt nhất bằng cách đánh giá \textbf{chi phí} tách; tách tốt nhất trở thành nút trong cây.
\end{phuongphap}

\subsubsection{Xây dựng cây (Building a Tree)}

\begin{ynghia}[title={Nút gốc và nút lá}]
Cây có nút gốc (root) và nút lá (terminal nodes).
Sau khi tạo nút gốc, xây cây theo hai phần sau.
\end{ynghia}

\begin{phuongphap}[title={Phần 1: Tạo nút lá (terminal node)}]
Quyết định khi nào \textbf{dừng} mở rộng cây, dùng hai tiêu chí:
\begin{itemize}
  \item \textbf{Maximum Tree Depth}: số nút tối đa sau nút gốc; dừng thêm nút lá khi đạt độ sâu tối đa.
  \item \textbf{Minimum Node Records}: số mẫu huấn luyện tối thiểu một nút chịu trách nhiệm; dừng khi đạt hoặc dưới ngưỡng này.
\end{itemize}
Nút lá dùng để đưa ra \textbf{dự đoán cuối cùng}.
\end{phuongphap}

\begin{phuongphap}[title={Phần 2: Tách đệ quy (Recursive Splitting)}]
Sau khi tạo một nút, tạo nút con \textbf{đệ quy} trên từng nhóm dữ liệu sinh ra bởi tách --- gọi lại cùng hàm nhiều lần.
\end{phuongphap}

\subsection{Dự đoán (Prediction)}

\begin{phuongphap}[title={Điều hướng trên cây}]
Sau khi xây cây, dự đoán bằng cách \textbf{điều hướng} cây với một hàng dữ liệu cụ thể.
Dùng hàm đệ quy tương tự khi xây cây: gọi lại routine dự đoán với nút con trái hoặc phải.
\end{phuongphap}

\subsection{Giả định}

\begin{luuy}[title={Giả định khi tạo decision tree}]
\begin{itemize}
  \item Tập huấn luyện là \textbf{nút gốc}.
  \item Bộ phân loại cây quyết định \textbf{ưu tiên} feature dạng phân loại (categorical). Nếu dùng giá trị liên tục cần \textbf{rời rạc hóa} (discretize) trước khi xây mô hình.
  \item Bản ghi được phân phối \textbf{đệ quy} theo giá trị thuộc tính.
  \item Dùng \textbf{phương pháp thống kê} để đặt thuộc tính ở vị trí nút (gốc hoặc nút trong).
\end{itemize}
\end{luuy}

\subsection{Triển khai bằng Python}

\begin{dongco}[title={Dataset Iris}]
150 mẫu hoa iris, bốn feature: chiều dài/rộng đài, chiều dài/rộng cánh hoa.
Ba lớp: setosa, versicolor, virginica.
\end{dongco}

\begin{vidu}[title={Import, nạp dữ liệu và chia tập}]
\begin{lstlisting}
import numpy as np
from sklearn.datasets import load_iris
from sklearn.model_selection import train_test_split
from sklearn.tree import DecisionTreeClassifier

iris = load_iris()
X_train, X_test, y_train, y_test = train_test_split(
    iris.data, iris.target, test_size=0.3, random_state=0)
\end{lstlisting}
\end{vidu}

\begin{vidu}[title={Huấn luyện và dự đoán}]
\begin{lstlisting}
dtc = DecisionTreeClassifier()
dtc.fit(X_train, y_train)
y_pred = dtc.predict(X_test)
\end{lstlisting}
\end{vidu}

\begin{vidu}[title={Độ chính xác}]
\begin{lstlisting}
accuracy = np.sum(y_pred == y_test) / len(y_test)
print("Accuracy:", accuracy)
\end{lstlisting}

\textbf{Output:}
\begin{verbatim}
Accuracy: 0.9777777777777777
\end{verbatim}
\end{vidu}

\begin{vidu}[title={Trực quan hóa cây}]
\begin{lstlisting}
import matplotlib.pyplot as plt
from sklearn.tree import plot_tree

plt.figure(figsize=(20, 10))
plot_tree(dtc, filled=True,
          feature_names=iris.feature_names,
          class_names=iris.target_names)
plt.show()
\end{lstlisting}

Hàm \texttt{plot\_tree} vẽ cấu trúc cây: mỗi nút là quyết định theo feature; mỗi lá là lớp hoặc giá trị số.
Tham số \texttt{filled=True} tô màu nút; \texttt{feature\_names} và \texttt{class\_names} gắn nhãn.
\end{vidu}

\begin{ynghia}[title={Biểu đồ cây quyết định}]
Biểu đồ minh họa:

\tsimg{05_decision_tree/img02.jpg}

Màu núc thể hiện lớp/ giá trị chiếm đa số trong nút; số ở đáy là số mẫu đi tới nút đó.
\end{ynghia}

\subsection{Ví dụ triển khai đầy đủ}

\begin{vidu}[title={Mã hoàn chỉnh}]
\begin{lstlisting}
import numpy as np
from sklearn.datasets import load_iris
from sklearn.model_selection import train_test_split
from sklearn.tree import DecisionTreeClassifier
import matplotlib.pyplot as plt
from sklearn.tree import plot_tree

iris = load_iris()
X_train, X_test, y_train, y_test = train_test_split(
    iris.data, iris.target, test_size=0.3, random_state=0)

dtc = DecisionTreeClassifier()
dtc.fit(X_train, y_train)
y_pred = dtc.predict(X_test)

accuracy = np.sum(y_pred == y_test) / len(y_test)
print("Accuracy:", accuracy)

plt.figure(figsize=(20, 10))
plot_tree(dtc, filled=True,
          feature_names=iris.feature_names,
          class_names=iris.target_names)
plt.show()
\end{lstlisting}

\textbf{Output:} \texttt{Accuracy: 0.9777777777777777} (khoảng \textbf{97,8\%}).
\end{vidu}
